\documentclass[a4paper,10pt,twoside]{article}
\usepackage{alefilot}
\begin{document}
\pagenumbering{gobble}

% ==============================================================================
%	NAME AND CONTACT INFORMATION
% ==============================================================================
\par{\centering{\Huge Δρ.\ ΦΙΛΟΘΕΟΥ, Αλέξανδρος}\bigskip\par}

\begin{center}
\vspace{-0.5em}
  {\small
   Λέβεν, Βέλγιο $\cdot$
  \href{https://github.com/li9i/}{\texttt{github.com/li9i}} $\ \cdot\ $ \href{https://scholar.google.com/citations?hl=en&view_op=list_works&gmla=AKzYXQ2ZNs_wAOy2e75F24-NJ0AGC7eeo8hugORYe3HG61yfKTR1H8vsWVf-C05DUq2qxvo9xqCvr6kF3qWTB4dIM-IJ&user=9_hI4hMAAAAJ}{google.scholar} $\ \cdot\ $ \href{https://docs.google.com/viewerng/viewer?url=https://raw.githubusercontent.com/li9i/portfolio/master/portfolio.pdf}{Portfolio} $\ \cdot \ $ \href{https://docs.google.com/viewerng/viewer?url=https://raw.githubusercontent.com/li9i/cv/master/references/references_alexandros_filotheou.pdf}{Συστάσεις} \\
  \href{mailto:alexandros.filotheou@gmail.com}{alexandros.filotheou@gmail.com} $\cdot$ (+30) 693 8787 677 $\cdot$ \href{https://www.linkedin.com/in/alexandros-filotheou/}{linkedin.com/alexandros-filotheou}
  }
\end{center}\vspace{-0.75em}

% ==============================================================================
%	SUMMARY
% ==============================================================================
\begin{gg_box} \small
Μηχανικός Ρομποτικής με οκτώ χρόνια σχεδίασης και εγκατάστασης εφαρμογών αυτόνομης πλοήγησης, εκτίμησης κατάστασης και οργάνωσης πολλαπλών ρομπότ, σε επίγειες και εναέριες πλατφόρμες, σε πραγματικές συνθήκες. Αποδεδειγμένη εμπειρία σε σύνθετες ενσωματώσεις πολλών εταίρων, από την αρχιτεκτονική έως την υλοποίηση και την εγκατάσταση στο πεδίο, σε πολλές χώρες. Με δημοσιεύσεις IEEE και μονοσυγγραφικά άρθρα στα IROS'24 και R\&AS.
\end{gg_box}

% ==============================================================================
% SKILLS
% ==============================================================================
\section{\textbf{Δεξιότητες}}\verticalspacebetweensectionandcontent\vspace{-0.75em}

\begin{center}
\begin{tabular}{rl}
Αγγλικά & Άριστα (IELTS 8.5 - C2 Proficiency) \\
Γλώσσες & \texttt{C/C++}, \texttt{Python}, \texttt{shell}, \texttt{MATLAB/Octave} \\
Ρομποτική / ΛΣ & \texttt{Linux},  \texttt{ROS}/\texttt{ROS 2} \\
Εργαλεία / Πλαίσια & \texttt{git}, \texttt{Docker}, Δέντρα Συμπεριφοράς, Gazebo, \texttt{Eigen}, \texttt{CI}/\texttt{CD}, \texttt{Qt}/\texttt{Tkinter} \\
Αισθητήρες / Υλικό & {\small LiDAR Ouster 3D και Hokuyo 2D, SICK Visionary-T, Astra, Realsense, ZED2, Impinj R420 }%\\
%SLAM, Πλοήγηση/Έλεγχος & \texttt{amcl}, \texttt{slam\_toolbox}, \texttt{rtabmap}, \texttt{openkarto}, \texttt{teb\_local\_planner}, \texttt{nav2}, MPC
\end{tabular}
\end{center}
\verticalspacebetweensections

% ==============================================================================
% ACHIEVEMENTS
% ==============================================================================
\section{\textbf{Επιτεύγματα}}\verticalspacebetweensectionandcontent
  \noindent\begin{minipage}[t]{\textwidth}
    \begin{itemize}
      \item Εγκατάσταση του ρομποτικού έργου κατασκευών \href{https://www.robetarme-project.eu/sobre/}{RoBétArmé} (\texteuro10Μ) στο \href{https://www.gem.wiki/Hermillon_hydroelectric_plant}{υδροηλεκτρικό εργοστάσιο Hermillon της EDF} με την Bouygues Construction, μετά από 2 χρόνια ενσωμάτωσης μεταξύ 12 εταίρων\verticalspaceafteritem
      \item Σχεδίαση και εγκατάσταση \href{https://www.youtube.com/watch?v=2EvTGNOqTrs}{αυτόνομου κοινωνικού ρομπότ} (NLP, πλοήγηση, SLAM, RFID, διεπαφή) στο \href{https://www.amth.gr/}{Αρχαιολογικό Μουσείο Θεσσαλονίκης}, $2$$.$$000$+ επισκέπτες ετησίως από το $2023$\verticalspaceafteritem
      \item Εντοπισμός συμφορήσεων $\rightarrow$ μείωση χρόνου εγκατάστασης >$10$x σε ρομπότ προσθετικής κατασκευής\verticalspaceafteritem
      \item Αύξηση της ακρίβειας εντοπισμού κατά $50\%$ σε στόλο ρομπότ απογραφής RFID $\rightarrow$ ακρίβεια \texttt{cm}\verticalspaceafteritem
      \item Επίλυση διαλειτουργικότητας \texttt{ROS}/\texttt{ROS~2} μεταξύ μηχανημάτων της κοινοπραξίας $\rightarrow$ αποδέσμευση και των 12 ομάδων\verticalspaceafteritem
      \item Εις εκ των δεκαοκτώ συγγραφέων μονοσυγγραφικών ανακοινώσεων στο \texttt{IEEE IROS'}$24$ ($1$$.$$587$ συνολικά)
    \end{itemize}
  \end{minipage} \\%[0.2em]
\verticalspacebetweensections

% ==============================================================================
%	LINKS - REPRESENTATIVE WORK
% ==============================================================================
\section{\textbf{Ενδεικτικό Έργο}}\verticalspacebetweensectionandcontent\vspace{-0.3em}

% \linkboxmany{fill}{frame}{corners}{hpad}{vpad}{text color}{content}

\begin{center}
\linkboxmany{blue!6}{blue!40}{3pt}{6pt}{1pt}{blue!70}{%
  \begin{tabular}{@{}c@{}}
    \textbf{Επιδείξεις} \\[-0.8em]
    \color{blue!40}\rule{4em}{0.4pt} \\[-0.2em]
    Ενσωμάτωση \lbref{https://www.youtube.com/shorts/TSIXd5L-Hxk}{[1]} $\cdot$
    \lbref{https://www.youtube.com/watch?v=ViU-9DLi2-M}{[2]} $\cdot$
    \lbref{https://www.youtube.com/watch?v=xuIbOL-MPRk\&t=10s}{[3]}\\
    \lbref{https://youtu.be/bo4lMI640DY?t=10}{SLAM/Πλοήγηση} \\
    \lbref{https://cultureid.web.auth.gr/?page\_id=200\&lang=en}{Ρομπότ Μουσείου} \\
    \lbref{https://www.youtube.com/watch?v=937OZez1iN8}{Control is Fun}
  \end{tabular}%
}
\quad
\linkboxmany{teal!10}{teal!50}{3pt}{6pt}{1pt}{teal!70}{%
  \begin{tabular}{@{}c@{}}
    \textbf{Εκτίμηση} \\
    \textbf{Κατάστασης} \\[-0.8em]
    \color{teal!50}\rule{6em}{0.4pt} \\[-0.2em]
    \lbref{https://www.youtube.com/watch?v=xaDKjI0WkDc\&t=11s}{Βίντεο} \\
    \lbref{https://ieeexplore.ieee.org/abstract/document/10802235}{Άρθρο} \\
    \lbref{https://github.com/cbgl/cbgl}{Κώδικας}
  \end{tabular}%
}
\quad
\linkboxmany{violet!10}{violet!40}{3pt}{6pt}{1pt}{violet!70}{%
  \begin{tabular}{@{}c@{}}
    \textbf{Οδομετρία} \\
    \textbf{Lidar} \\[-0.8em]
    \color{violet!40}\rule{6em}{0.4pt} \\[-0.2em]
    \lbref{https://www.youtube.com/watch?v=hB4qsHCEXGI\&t=11s}{Βίντεο} \\
    \lbref{https://ieeexplore.ieee.org/abstract/document/9981228}{Άρθρο} \\
    \lbref{https://github.com/fourier-scan-matcher/fsm-lo}{Κώδικας}
  \end{tabular}%
}\quad
\linkboxmany{yellow!20}{orange!50}{3pt}{6pt}{1pt}{orange!70}{%
  \begin{tabular}{@{}c@{}}
    \textbf{Δημοσιεύσεις} \\[-0.8em]
    \color{orange!50}\rule{7em}{0.4pt} \\[-0.2em]
    Εκτίμηση \lbref{https://ieeexplore.ieee.org/abstract/document/9981228}{[Τοπική]} $\cdot$ \lbref{https://ieeexplore.ieee.org/abstract/document/10802235}{[Ολική]} \\
    \lbref{https://www.sciencedirect.com/science/article/abs/pii/S0921889021002323}{Εντοπισμός/Έλεγχος (μόνος)} \\
    \lbref{https://www.tandfonline.com/doi/full/10.1080/00207179.2018.1514129}{MPC πολλαπλών ρομπότ} \\
    \lbref{https://link.springer.com/article/10.1007/s10846-019-01086-y}{Επισκόπηση Πλοήγησης}
  \end{tabular}%
}

\end{center}
\verticalspacebetweensections\vspace{-0.25em}

% ==============================================================================
%	WORK EXPERIENCE
% ==============================================================================
\section{\textbf{Εργασιακή Εμπειρία}}\verticalspacebetweensectionandcontent\vspace{+0.25em}

  \noindent\textbf{Senior Robotics Software Engineer} $\cdot$ \href{https://www.intermodalics.ai/}{Intermodalics}, Λέβεν ΒΕ\hfill {\small \textcolor{agray}{Φεβ $2026$ -- σήμερα}} \\
  \begin{minipage}[t]{\textwidth}
    \begin{itemize}
      \item Για την \href{https://www.youtube.com/watch?v=LWqDSMOAeec&t=10s}{KION Group AG}: διασύνδεση διαχειριστή στόλου VDA5050 με συμπεριφορές αποστολών \texttt{ROS~2}, σε διεθνή ομάδα ανάπτυξης αυτόνομων περονοφόρων AMR μαζικής παραγωγής για αποθήκες
    \end{itemize}
  \end{minipage} \\[-0.2em]

  \noindent\textbf{Μηχανικός Ενσωμάτωσης \& Εγκατάστασης} $\cdot$ \href{https://robotics.iti.gr/en/home/}{ΕΚΕΤΑ/ΙΠΤΗΛ}, Θεσσαλονίκη ΕΛ \hfill {\small \textcolor{agray}{Σεπ $2023$ -- Ιαν $2026$}} \\
  \begin{minipage}[t]{\textwidth}
    \begin{itemize}
      \item Μοναδικός αρχιτέκτονας και ενσωματωτής λογισμικού του έργου Horizon \href{https://www.robetarme-project.eu/sobre/}{RoBétArmé} (GA 101058731)\verticalspaceafteritem
    \item Καθοδήγηση σχεδιασμού εργασιών και δέντρων συμπεριφοράς για ρομπότ προσθετικής κατασκευής σκυροδέματος και \href{https://www.youtube.com/watch?v=xuIbOL-MPRk}{μετάλλου}\verticalspaceafteritem
    \item Καθορισμός αρχών ανάπτυξης, ενσωμάτωσης, εγκατάστασης για $12$ ομάδες και >$50$ πακέτα \texttt{ROS~2} σε Docker
    \end{itemize}
  \end{minipage} \\[-0.2em]

  \noindent\textbf{Εργολήπτης Ερευνητικών Έργων Ρομποτικής} $\cdot$ Τμ. ΗΜΜΥ, Α.Π.Θ., Θεσσαλονίκη ΕΛ \hfill {\small \textcolor{agray}{Σεπ $2018$ -- Μαρ $2023$}} \\
  \begin{minipage}[t]{\textwidth}
    \begin{itemize}
      \item Μοναδικός προγραμματιστής του ρομποτικού λογισμικού των έργων \href{https://relief.web.auth.gr/language/en/home/\#pg-72-1}{RELIEF} και \href{https://cultureid.web.auth.gr/?page\_id=200&lang=en}{CultureId}: επίγειες (Turtlebot~2, Robotnik RB1) και εναέριες πλατφόρμες απογραφής RFID, κοινωνικό ρομπότ ξενάγησης μουσείου\verticalspaceafteritem
      \item Παράδοση 2D/3D SLAM και αυτόνομης πλοήγησης σε αποθήκες, βιβλιοθήκες, μουσεία \verticalspaceafteritem
      \item Ανάπτυξη και ενσωμάτωση κώδικα που τροφοδότησε $18$+ δημοσιεύσεις σε κορυφαία περιοδικά και συνέδρια IEEE
    \end{itemize}
  \end{minipage} \\[-0.2em]

\noindent \textbf{Μηχανικός Μηχανικής Όρασης} $\cdot$ \href{https://issel.ee.auth.gr/pandora-robotics/}{Ομάδα ρομποτικής PANDORA}, Θεσσαλονίκη ΕΛ \hfill {\small \textcolor{agray}{Οκτ $2013$ -- Ιουλ $2014$}}
  \begin{minipage}[t]{\textwidth}
    \begin{itemize}
      \item $2^{\textnormal{η}}$ θέση στον RoboCup Rescue $2015$ με \href{https://github.com/li9i/pandora\_vision\_2014}{συστήματος εύρεσης οπών σε τοίχους σε C++} με Kinect
    \end{itemize}
  \end{minipage} \\ [-0.8em]

\verticalspacebetweensections

% ==============================================================================
%	VOLUNTARY EXPERIENCE
% ==============================================================================
\section{\textbf{Συνεισφορές σε Λογισμικό Ανοικτού Κώδικα}}\verticalspacebetweensectionandcontent

  \noindent\begin{minipage}[t]{\textwidth}
    \begin{itemize}
      \item (Linux) \href{https://github.com/li9i/gitrl-z}{\texttt{gitrl-z}} $\cdot$ Εργαλείο πλοήγησης στο reflog του git για επαναφορά του δέντρου μετά από λάθη\verticalspaceafteritem
      \item (ROS 2) \href{https://index.ros.org/p/hitch_estimation_apriltag_array/#humble}{\texttt{hitch\_estimation\_apriltag\_array}} $\cdot$ Εκτίμηση γωνίας μεταξύ οχήματος και ρυμουλκούμενου \verticalspaceafteritem
      \item (ROS 2) \href{https://index.ros.org/r/pointcloud_conversions/#lyrical}{\texttt{pointcloud\_conversions}} $\cdot$ Μετατροπές μορφής και συναρτήσεις αποθήκευσης για νέφη σημείων
    \end{itemize}
  \end{minipage} \\ [0.2em]

\verticalspacebetweensections

% ==============================================================================
%	EDUCATION
% ==============================================================================
\section{\textbf{Εκπαίδευση}}%\verticalspacebetweensectionandcontent
\vspace{-1em}
\begin{center}
\begin{tabular}{ll}
  \href{https://www.didaktorika.gr/eadd/handle/10442/56399?locale=el}{\textbf{Διδακτορικό}} $\cdot$ Ηλεκτρολόγων Μηχ. \& Μηχ. Υπολογιστών $\cdot$ Α.Π.Θ. & {\small \textcolor{agray}{Σεπ $2018$ -- Ιουν $2023$}} \\
  %& \hspace{-30.3em} Εκπονήθηκε και υποστηρίχθηκε χωρίς επιβλέποντα \\
  \href{http://kth.diva-portal.org/smash/record.jsf?pid=diva2\%3A1102597\&dswid=2875}{\textbf{Μεταπτυχιακό}} $\cdot$ Systems, Control, and Robotics $\cdot$ KTH Royal Inst. of Technology & {\small \textcolor{agray}{Σεπ $2015$ -- Ιουν $2017$}} \\
  & \hspace{-30.3em} Εργασία: Αποκεντρωμένος MPC για πλοήγηση πολλαπλών ρομπότ [ECC'18, IJC'18] \\
  \href{https://ikee.lib.auth.gr/record/291560}{\textbf{Δίπλωμα}} $\cdot$ Ηλεκτρολόγων Μηχ. \& Μηχ. Υπολογιστών $\cdot$ Α.Π.Θ. & {\small \textcolor{agray}{Σεπ $2005$ -- Ιουλ $2013$}}
\end{tabular}
\end{center}

\end{document}
